\documentclass[a4paper,10pt,brazil]{article}
\usepackage{provastyle}
\usepackage{menukeys}

\begin{document}
\makeexamtitlepage{UFOP}{BCC222: Programação Funcional}{José Romildo}{2026/1}{Primeira Prova}{03}
\begin{multicols}{2}
\questionTitle{1}{Inferência de tipo de isDigit}
\begin{flushright}
  \footnotesize
  \textit{\textbf{4: Tipos de Dados}}\\
  \textit{c04-inferencia-basica-02}\\

\end{flushright}
Considere \haskellinline{import Data.Char}. Qual é o tipo inferido para a expressão \haskellinline{isDigit '5'}?

\begin{enumerate}
  \item \haskellinline{Char}

  \item \haskellinline{Char -> Bool}

  \item \haskellinline{Num a => a}

  \item \haskellinline{True}

  \item \haskellinline{Bool}


\end{enumerate}

\questionTitle{2}{Falha de readLn diante de entrada inválida}
\begin{flushright}
  \footnotesize
  \textit{\textbf{10: Ações de E/S recursivas}}\\
  \textit{c10-bordas-robustez-01}\\

\end{flushright}
No capítulo, os exemplos usam \haskellinline|readLn| para simplificar a leitura. O que acontece, em geral, se uma ação espera um número e o usuário digita um texto que não pode ser convertido para o tipo esperado?

\begin{enumerate}
  \item O valor inválido é interpretado como zero, pois zero é o valor neutro das leituras numéricas.

  \item \haskellinline|readLn| devolve automaticamente \haskellinline|Nothing|, pois seu tipo padrão é \haskellinline|IO (Maybe a)|.

  \item \haskellinline|readLn| repete a pergunta automaticamente até receber uma entrada válida.

  \item O compilador detecta a entrada inválida antes da execução do programa.

  \item A ação pode falhar em tempo de execução; uma versão mais robusta deveria validar a entrada, por exemplo com uma estratégia baseada em \haskellinline|readMaybe|.


\end{enumerate}

\questionTitle{3}{Trabalho pendente após a chamada recursiva}
\begin{flushright}
  \footnotesize
  \textit{\textbf{10: Ações de E/S recursivas}}\\
  \textit{c10-acumuladores-cauda-01}\\

\end{flushright}
Considere a versão sem acumulador:

\begin{haskellcode}
lerESomar :: IO Integer
lerESomar = do
  n <- readLn
  if n == 0
    then return 0
    else do
      somaResto <- lerESomar
      return (n + somaResto)
\end{haskellcode}

Por que a chamada recursiva não está em posição de cauda?

\begin{enumerate}
  \item Porque, depois que \haskellinline|lerESomar| retorna, ainda é necessário calcular \haskellinline|n + somaResto| e aplicar \haskellinline|return| ao resultado.

  \item Porque o tipo \haskellinline|IO Integer| é incompatível com recursão de cauda.

  \item Porque \haskellinline|if| não permite chamadas recursivas em posição de cauda.

  \item Porque o caso base usa \haskellinline|return 0| em vez de \haskellinline|return n|.

  \item Porque qualquer uso de \haskellinline|readLn| impede uma função de ser recursiva de cauda.


\end{enumerate}

\questionTitle{4}{Interpretação de Num a => a}
\begin{flushright}
  \footnotesize
  \textit{\textbf{4: Tipos de Dados}}\\
  \textit{c04-ghci-interpretacao-01}\\

\end{flushright}
Ao executar \texttt{:type 42} no GHCi, a resposta é \haskellinline{42 :: Num a => a}. O que isso significa?

\begin{enumerate}
  \item \haskellinline{42} é um valor de um tipo desconhecido que não pode ser usado em operações aritméticas.

  \item \haskellinline{42} é uma constante numérica que só pode ser do tipo \haskellinline{Int}.

  \item Significa que \haskellinline{42} é uma função que recebe um \haskellinline{Num} e retorna um \haskellinline{a}.

  \item \haskellinline{42} pode ser interpretado como um valor de qualquer tipo que pertença à classe \haskellinline{Num} (como \haskellinline{Int}, \haskellinline{Integer}, \haskellinline{Double}).

  \item Indica um erro de tipo, pois o tipo é ambíguo e não pode ser resolvido.


\end{enumerate}

\questionTitle{5}{Equivalência com True}
\begin{flushright}
  \footnotesize
  \textit{\textbf{5: Expressão condicional}}\\
  \textit{c05-otherwise-02}\\

\end{flushright}
Qual definição é equivalente ao uso usual de \haskellinline|otherwise| na última guarda?

\begin{enumerate}
  \item Substituir \haskellinline|otherwise| por \haskellinline|True|.

  \item Substituir \haskellinline|otherwise| por \haskellinline|undefined|.

  \item Substituir \haskellinline|otherwise| por \haskellinline|False|.

  \item Substituir \haskellinline|otherwise| por \haskellinline|else|.

  \item Substituir \haskellinline|otherwise| por \haskellinline|()|.


\end{enumerate}

\questionTitle{6}{O significado de IO ()}
\begin{flushright}
  \footnotesize
  \textit{\textbf{8: Programas interativos}}\\
  \textit{c08-tipos-acoes-01}\\

\end{flushright}
Por que muitas ações de saída, como \haskellinline|putStrLn :: String -> IO ()|, produzem um valor do tipo \haskellinline|IO ()|?

\begin{enumerate}
  \item Porque a ação realiza um efeito de E/S, mas não produz nenhum dado significativo para ser usado pelo restante do programa. O valor produzido é o valor unitário \haskellinline|()|.

  \item Porque \haskellinline|IO ()| é um sinônimo de \haskellinline|IO String|.

  \item Porque o tipo \haskellinline|()| indica que a ação pode ser executada no máximo uma vez.

  \item Porque a função não recebe argumentos.

  \item Porque a ação sempre falha e retorna um valor vazio.


\end{enumerate}

\questionTitle{7}{Lookup com valor padrão}
\begin{flushright}
  \footnotesize
  \textit{\textbf{6: Estruturas de dados básicas}}\\
  \textit{c06-lookup-associacao-02}\\

\end{flushright}
Qual é o resultado da expressão abaixo?
\begin{minted}{haskell}
let tabela = [("ana",8.5),("bruno",7.0)]
in fromMaybe 0.0 (lookup "davi" tabela)
\end{minted}

\begin{enumerate}
  \item \haskellinline{ Nothing }

  \item \haskellinline{ 7.0 }

  \item \haskellinline{ Just 0.0 }

  \item Ocorre um erro em tempo de execução.

  \item \haskellinline{ 0.0 }


\end{enumerate}

\questionTitle{8}{A função putStr}
\begin{flushright}
  \footnotesize
  \textit{\textbf{8: Programas interativos}}\\
  \textit{c08-entrada-saida-01}\\

\end{flushright}
Qual é o tipo de \haskellinline|putStr| e qual é seu comportamento?

\begin{enumerate}
  \item \haskellinline|String -> IO ()|. Ela constrói uma ação que imprime a string na saída padrão sem acrescentar uma quebra de linha.

  \item \haskellinline|String -> String|. Ela apenas converte uma string em outra string.

  \item \haskellinline|IO String|. Ela lê uma string da entrada padrão.

  \item \haskellinline|String -> IO String|. Ela imprime a string e também a retorna para o programa.

  \item \haskellinline|IO ()|. Ela imprime uma linha vazia sem receber argumentos.


\end{enumerate}

\questionTitle{9}{Construção da lista durante o retorno da recursão}
\begin{flushright}
  \footnotesize
  \textit{\textbf{10: Ações de E/S recursivas}}\\
  \textit{c10-listas-ordem-02}\\

\end{flushright}
Analise a versão sem acumulador:

\begin{haskellcode}
lerLista :: IO [Integer]
lerLista = do
  x <- readLn
  if x == 0
    then return []
    else do
      resto <- lerLista
      return (x:resto)
\end{haskellcode}

Se o usuário digita \haskellinline|4|, depois \haskellinline|9|, e depois \haskellinline|0|, qual lista é produzida?

\begin{enumerate}
  \item \haskellinline|[4,9]|, pois cada chamada adiciona seu valor à frente da lista produzida pelo restante da recursão.

  \item \haskellinline|[0,9,4]|, pois a recursão reconstrói a lista de trás para frente incluindo o caso base.

  \item \haskellinline|[]|, pois \haskellinline|return []| descarta os valores anteriores.

  \item \haskellinline|[9,4]|, pois cada novo elemento é colocado no início da lista final.

  \item \haskellinline|[4,9,0]|, pois o sentinela também faz parte da lista retornada.


\end{enumerate}

\questionTitle{10}{If como subexpressão}
\begin{flushright}
  \footnotesize
  \textit{\textbf{5: Expressão condicional}}\\
  \textit{c05-if-avaliacao-02}\\

\end{flushright}
Considere a expressão:
\begin{haskellcode}
4 * (if 'B' < 'A' then 2 + 3 else 2) - 3
\end{haskellcode}
Qual é o seu valor?

\begin{enumerate}
  \item \haskellinline|5|

  \item \haskellinline|17|

  \item Ocorre um erro de tipos.

  \item \haskellinline|-1|

  \item \haskellinline|20|


\end{enumerate}

\questionTitle{11}{Estrutura geral de um laço recursivo de E/S}
\begin{flushright}
  \footnotesize
  \textit{\textbf{10: Ações de E/S recursivas}}\\
  \textit{c10-lacos-io-01}\\

\end{flushright}
No estilo usado no capítulo, qual é o padrão básico para repetir uma ação de E/S, como ler vários valores do teclado?

\begin{enumerate}
  \item Importar uma biblioteca obrigatória de laços, pois Haskell não permite repetição de ações de E/S sem ela.

  \item Definir uma ação recursiva que executa uma parte da interação e chama a si mesma até atingir um caso base.

  \item Usar \haskellinline|map| sobre uma lista infinita e confiar que o \haskellinline|IO| executará apenas os elementos necessários.

  \item Usar uma compreensão de listas dentro de \haskellinline|do| para repetir efeitos colaterais.

  \item Declarar uma variável mutável global e atualizá-la até que a condição de parada seja falsa.


\end{enumerate}

\questionTitle{12}{Resultados de funções RealFrac}
\begin{flushright}
  \footnotesize
  \textit{\textbf{7: Polimorfismo}}\\
  \textit{cap07-conversao-realfrac-param}\\

\end{flushright}
Qual é o resultado da avaliação da expressão abaixo?

\begin{minted}{haskell}
truncate (-2.8) :: Int
\end{minted}

\begin{enumerate}
  \item A expressão causa erro de tipo, pois \haskellinline{RealFrac} não permite retornar \haskellinline{Int}.
  \item \haskellinline|-2|
  \item \haskellinline|-3|
  \item \haskellinline|2|
  \item \haskellinline|3|

\end{enumerate}

\questionTitle{13}{Cabeça, cauda e construção}
\begin{flushright}
  \footnotesize
  \textit{\textbf{6: Estruturas de dados básicas}}\\
  \textit{c06-listas-estrutura-02}\\

\end{flushright}
Qual é o resultado da expressão abaixo?
\begin{minted}{haskell}
head (tail [1,2,3,4])
\end{minted}

\begin{enumerate}
  \item \haskellinline{ 1 }

  \item \haskellinline{ 2 }

  \item \haskellinline{ 3 }

  \item A expressão resulta em erro de execução.

  \item \haskellinline{ [2,3,4] }


\end{enumerate}

\questionTitle{14}{Paradigmas - Princípio Nuclear (Legado Refatorado)}
\begin{flushright}
  \footnotesize
  \textit{\textbf{1: Paradigmas de programação}}\\
  \textit{c01-conceitos-paradigmas-02}\\

\end{flushright}
Ao contrastarmos os paradigmas dominantes, o \textbf{princípio nuclear} que separa a programação imperativa da programação declarativa (funcional) reside no foco da estruturação do código. Qual é essa diferença?

\begin{enumerate}
  \item A programação imperativa exige a declaração estrita de tipos em tempo de compilação, enquanto a declarativa permite que os tipos sejam inferidos dinamicamente durante a execução.

  \item A programação imperativa é estruturada em torno de funções matemáticas puras, enquanto a declarativa é estruturada em torno de objetos que trocam mensagens.

  \item A programação imperativa abstrai o gerenciamento de memória através de \emph{Garbage Collectors}, enquanto a declarativa exige alocação explícita com \cinline|malloc|.

  \item A programação imperativa foca em descrever \emph{como} o computador deve executar os passos (controle de fluxo e estado), enquanto a declarativa foca em descrever \emph{o que} é o resultado desejado (lógica e propriedades).

  \item A programação imperativa foi projetada exclusivamente para arquiteturas mononucleares (single-core), enquanto a declarativa foi projetada nativamente para paralelismo em GPUs.


\end{enumerate}

\questionTitle{15}{Modelo de E/S em uma linguagem pura}
\begin{flushright}
  \footnotesize
  \textit{\textbf{8: Programas interativos}}\\
  \textit{c08-modelo-io-01}\\

\end{flushright}
Como Haskell consegue realizar operações de entrada e saída e, ainda assim, preservar a ideia de linguagem puramente funcional?

\begin{enumerate}
  \item Valores do tipo \haskellinline|IO a| são descrições de ações interativas. O programa constrói essas descrições de forma pura, e o \emph{runtime system} é responsável por executá-las.

  \item Haskell não realiza entrada e saída; todo programa interativo precisa ser escrito em outra linguagem e apenas chamado a partir de Haskell.

  \item As funções de entrada e saída são exceções impuras dentro da linguagem e, por isso, não participam do sistema de tipos.

  \item O tipo \haskellinline|IO a| armazena internamente uma variável global mutável que pode ser alterada por qualquer função pura.

  \item A notação \haskellinline|do| transforma temporariamente Haskell em uma linguagem imperativa, desativando a transparência referencial.


\end{enumerate}

\questionTitle{16}{Otimização de chamada de cauda em Haskell}
\begin{flushright}
  \footnotesize
  \textit{\textbf{9: Funções recursivas}}\\
  \textit{c09-eficiencia-tco-01}\\

\end{flushright}
Qual afirmação é a mais adequada sobre recursividade de cauda e uso de memória em Haskell?

\begin{enumerate}
  \item Recursividade de cauda ajuda a evitar operações pendentes, mas em Haskell o uso de memória também depende de como os acumuladores são avaliados.

  \item Recursividade de cauda é impossível em Haskell, pois a linguagem é puramente funcional.

  \item O uso de acumuladores sempre elimina a necessidade de casos base.

  \item Toda função recursiva de cauda em Haskell usa memória constante, independentemente da avaliação dos argumentos.

  \item Haskell rejeita qualquer função que não seja recursiva de cauda.


\end{enumerate}

\questionTitle{17}{Saudação interativa}
\begin{flushright}
  \footnotesize
  \textit{\textbf{8: Programas interativos}}\\
  \textit{c08-programas-completos-01}\\

\end{flushright}
Considere o programa:

\begin{minted}{haskell}
main :: IO ()
main = do
  putStrLn "Nome?"
  nome <- getLine
  putStrLn ("Olá, " ++ nome ++ "!")
\end{minted}

Se o usuário digitar \haskellinline|Ana|, qual será a última linha impressa?

\begin{enumerate}
  \item \haskellinline|Olá, nome!|

  \item \haskellinline|Olá, getLine!|

  \item \haskellinline|Ana|

  \item \haskellinline|Nome?|

  \item \haskellinline|Olá, Ana!|


\end{enumerate}

\questionTitle{18}{let e <- em blocos do}
\begin{flushright}
  \footnotesize
  \textit{\textbf{8: Programas interativos}}\\
  \textit{c08-blocos-do-01}\\

\end{flushright}
Dentro de um bloco \haskellinline|do| de uma ação \haskellinline|IO|, qual é a diferença fundamental entre \haskellinline|x <- acao| e \haskellinline|let x = expr|?

\begin{enumerate}
  \item \haskellinline|x <- acao| executa uma ação de E/S e associa a \haskellinline|x| o valor produzido por ela; \haskellinline|let x = expr| define localmente um valor ou função pura, sem executar ação de E/S.

  \item \haskellinline|<-| só pode ser usado para valores numéricos; \haskellinline|let| só pode ser usado para strings.

  \item \haskellinline|let| executa a ação de E/S, enquanto \haskellinline|<-| apenas cria um apelido para a ação.

  \item Não há diferença semântica; as duas formas são intercambiáveis em blocos \haskellinline|do|.

  \item \haskellinline|let| só é permitido fora de blocos \haskellinline|do|.


\end{enumerate}

\questionTitle{19}{Anatomia de uma definição em arquivo}
\begin{flushright}
  \footnotesize
  \textit{\textbf{3: Expressões e definições}}\\
  \textit{c03-arquivos-identificadores-02}\\

\end{flushright}
Considere a definição em um arquivo fonte:

\begin{haskellcode}
dobro n = 2 * n
\end{haskellcode}

Qual alternativa identifica corretamente seus componentes?

\begin{enumerate}
  \item A definição é inválida porque funções de um argumento exigem vírgula depois do nome.

  \item \haskellinline|2 * n| é uma assinatura de tipo.

  \item \haskellinline|dobro| é o identificador definido, \haskellinline|n| é o parâmetro e \haskellinline|2 * n| é o corpo da definição.

  \item \haskellinline|dobro n| é o nome completo da função, e por isso ela não recebe argumentos.

  \item \haskellinline|n| é uma definição global independente criada automaticamente.


\end{enumerate}

\questionTitle{20}{Comentários aninhados e código efetivo}
\begin{flushright}
  \footnotesize
  \textit{\textbf{3: Expressões e definições}}\\
  \textit{c03-comentarios-valores-02}\\

\end{flushright}
Considere o trecho a seguir.

\begin{haskellcode}
x = 10
{-
y = 20
  {- z = 30 -}
w = 40
-}
r = x + 1
\end{haskellcode}

Quais definições permanecem como código efetivo para o compilador?

\begin{enumerate}
  \item Nenhuma definição, pois comentários de bloco aninhados invalidam o arquivo.

  \item \haskellinline|x|, \haskellinline|z| e \haskellinline|r|.

  \item Apenas \haskellinline|r|, pois \haskellinline|x| fica dentro do comentário de bloco.

  \item \haskellinline|x|, \haskellinline|y|, \haskellinline|w| e \haskellinline|r|.

  \item Apenas \haskellinline|x| e \haskellinline|r|.


\end{enumerate}

\questionTitle{21}{Declaração let no GHCi}
\begin{flushright}
  \footnotesize
  \textit{\textbf{3: Expressões e definições}}\\
  \textit{c03-ghci-definicoes-02}\\

\end{flushright}
No GHCi, qual alternativa representa a forma explícita de associar o nome \haskellinline|raio| ao valor \haskellinline|10| para uso em expressões posteriores?

\begin{enumerate}
  \item \haskellinline|raio <- 10|.

  \item \haskellinline|let 10 = raio|.

  \item \haskellinline|raio := 10|.

  \item \haskellinline|def raio = 10|.

  \item \haskellinline|let raio = 10|.


\end{enumerate}

\questionTitle{22}{Ineficiência do Fibonacci ingênuo}
\begin{flushright}
  \footnotesize
  \textit{\textbf{9: Funções recursivas}}\\
  \textit{c09-formas-fibonacci-01}\\

\end{flushright}
Considere a definição ingênua:
\begin{haskellcode}
fib n | n == 0    = 0
      | n == 1    = 1
      | otherwise = fib (n - 2) + fib (n - 1)
\end{haskellcode}
Por que essa definição é ineficiente para valores grandes de \haskellinline|n|?

\begin{enumerate}
  \item Porque ela armazena todos os resultados anteriores em uma lista global.

  \item Porque a chamada recursiva é de cauda e, portanto, nunca retorna.

  \item Porque \haskellinline|+| entre inteiros é proibido em funções recursivas.

  \item Porque ela recalcula os mesmos subproblemas muitas vezes, fazendo o número de chamadas crescer rapidamente.

  \item Porque a função não tem casos base para \haskellinline|0| e \haskellinline|1|.


\end{enumerate}

\questionTitle{23}{Assinatura de função com Bool e lista de Char}
\begin{flushright}
  \footnotesize
  \textit{\textbf{4: Tipos de Dados}}\\
  \textit{c04-funcao-assinatura-02}\\

\end{flushright}
Qual das seguintes assinaturas de tipo representa uma função que recebe um \haskellinline{Bool} e uma lista de \haskellinline{Char} e retorna um \haskellinline{Char}?

\begin{enumerate}
  \item \haskellinline{(Bool, [Char], Char)}

  \item \haskellinline{Bool -> [Char] -> Char}

  \item \haskellinline{Char -> Bool -> [Char]}

  \item \haskellinline{[Bool] -> [Char] -> Char}

  \item \haskellinline{(Bool, [Char]) -> Char}


\end{enumerate}

\questionTitle{24}{A Crise do Software e a solução funcional}
\begin{flushright}
  \footnotesize
  \textit{\textbf{1: Paradigmas de programação}}\\
  \textit{c01-engenharia-software-02}\\

\end{flushright}
O termo \emph{Crise do Software} surgiu historicamente devido à dificuldade de entregar sistemas complexos no prazo, dentro do orçamento e isentos de falhas catastróficas em tempo de execução. Como o paradigma funcional atua para mitigar diretamente a raiz estrutural dessa crise?

\begin{enumerate}
  \item Restringindo rigorosamente a arquitetura dos projetos a um limite máximo de linhas de código por pacote para evitar que os softwares atinjam uma escala insustentável.

  \item Elevando o nível de abstração focado na declaração (o que fazer) e adotando bases que eliminam o estado mutável não documentado, o que reduz drasticamente a complexidade do rastreamento de defeitos.

  \item Reduzindo o tempo de desenvolvimento ao automatizar a geração de código através de modelos probabilísticos baseados em inteligência artificial no sistema de tipos dinâmicos.

  \item Forçando a adoção de metodologias de desenvolvimento ágil (\emph{Agile/Scrum}) através de ferramentas nativas embutidas diretamente nos compiladores das linguagens funcionais.

  \item Eliminando de forma definitiva a necessidade de escrever casos de teste unitário, já que o processo de compilação cruzada garante a integridade incondicional das regras de negócio.


\end{enumerate}

\questionTitle{25}{Ordem de avaliação das guardas}
\begin{flushright}
  \footnotesize
  \textit{\textbf{5: Expressão condicional}}\\
  \textit{c05-guardas-sintaxe-semantica-02}\\

\end{flushright}
Como as guardas de uma definição são avaliadas?

\begin{enumerate}
  \item As guardas são avaliadas de baixo para cima.

  \item Todas as guardas são avaliadas, e a última verdadeira determina o resultado.

  \item O compilador reordena as guardas para testar primeiro as mais específicas.

  \item De cima para baixo; a primeira guarda verdadeira determina o resultado.

  \item O resultado é escolhido pela guarda cuja expressão à direita de \haskellinline|=| parece mais simples.


\end{enumerate}

\questionTitle{26}{Restrição Read e necessidade de contexto}
\begin{flushright}
  \footnotesize
  \textit{\textbf{7: Polimorfismo}}\\
  \textit{cap07-adhoc-read-contexto}\\

\end{flushright}
A função \haskellinline{read} tem assinatura:

\begin{minted}{haskell}
read :: Read a => String -> a
\end{minted}

Qual uso deixa claro para o compilador que o resultado deve ser um \haskellinline{Int}?

\begin{enumerate}
  \item \mintinline{haskell}{read "10" :: String -> Int}
  \item \mintinline{haskell}{show (read "10")}
  \item \mintinline{haskell}{read "10" :: Int}
  \item \mintinline{haskell}{read :: "10" -> Int}
  \item \mintinline{haskell}{length (read "10")}

\end{enumerate}

\questionTitle{27}{A saída segura do ambiente interativo}
\begin{flushright}
  \footnotesize
  \textit{\textbf{2: Ambiente de desenvolvimento}}\\
  \textit{c02-fluxo-trabalho-02}\\

\end{flushright}
Após terminar o desenvolvimento, testar exaustivamente o arquivo \texttt{.hs} e salvar o progresso, qual comando do GHCi é o responsável por encerrar a sessão interativa de forma limpa, devolvendo o controle ao terminal do sistema operacional?

\begin{enumerate}
  \item O comando \texttt{:close} (ou a abreviação \texttt{:c}).

  \item O comando \texttt{:quit} (ou a abreviação \texttt{:q}).

  \item O comando \texttt{:stop} (ou a abreviação \texttt{:s}).

  \item O comando \texttt{:end} (sem abreviação suportada).

  \item O comando \texttt{:exit} (ou a abreviação \texttt{:e}).


\end{enumerate}

\questionTitle{28}{Carregando módulos no GHCi}
\begin{flushright}
  \footnotesize
  \textit{\textbf{2: Ambiente de desenvolvimento}}\\
  \textit{c02-ghci-modulos-02}\\

\end{flushright}
Para utilizar funções avançadas de listas (como \haskellinline{sort}) diretamente no prompt do GHCi, o desenvolvedor precisa trazer o módulo \texttt{Data.List} para o escopo atual. Quais comandos são válidos para realizar esta operação?

\begin{enumerate}
  \item O comando \texttt{:browse Data.List}, que carrega e executa todas as funções contidas no módulo de forma simultânea.

  \item Não é necessário usar comandos; em Haskell, diferentemente de C ou Java, todos os módulos oficiais estão sempre injetados globalmente no Prelude.

  \item Apenas o comando interno \texttt{:load Data.List}, pois o termo \haskellinline{import} é exclusivo para arquivos físicos \texttt{.hs}.

  \item A palavra-chave nativa da linguagem \haskellinline{import Data.List} ou o atalho específico do GHCi \texttt{:module +Data.List} (ou \texttt{:m +Data.List}).

  \item O comando \texttt{:install Data.List}, que solicita a função diretamente do repositório remoto Hackage em tempo real.


\end{enumerate}

\questionTitle{29}{Associatividade da aplicação de função}
\begin{flushright}
  \footnotesize
  \textit{\textbf{3: Expressões e definições}}\\
  \textit{c03-aplicacao-notacoes-02}\\

\end{flushright}
A aplicação de função em Haskell associa à esquerda. Assim, como a expressão \haskellinline|logBase 2 8| deve ser lida?

\begin{enumerate}
  \item Como \haskellinline|logBase (2 8)|, pois o argumento mais à direita é agrupado primeiro.

  \item Como \haskellinline|(logBase 2 8)|, mas com \haskellinline|2| e \haskellinline|8| formando implicitamente uma tupla.

  \item Como \haskellinline|(logBase 2) 8|: primeiro obtém-se a função logarítmica de base \haskellinline|2|, depois ela é aplicada a \haskellinline|8|.

  \item Como \haskellinline|logBase (2, 8)|, pois funções de dois argumentos recebem tuplas por padrão.

  \item Como uma expressão inválida, pois \haskellinline|logBase| exige que os argumentos sejam separados por vírgula.


\end{enumerate}

\questionTitle{30}{Entrada que provoca falha}
\begin{flushright}
  \footnotesize
  \textit{\textbf{5: Expressão condicional}}\\
  \textit{c05-guardas-exaustividade-02}\\

\end{flushright}
Para qual chamada a função abaixo falha em tempo de execução?
\begin{haskellcode}
comparar x y
  | x > y = "Maior"
  | x < y = "Menor"
\end{haskellcode}

\begin{enumerate}
  \item \haskellinline|comparar 0 (-1)|

  \item \haskellinline|comparar 5 5|

  \item Nenhuma, pois o compilador adiciona \haskellinline|otherwise| implicitamente.

  \item \haskellinline|comparar 10 5|

  \item \haskellinline|comparar 5 10|


\end{enumerate}

\questionTitle{31}{Strings como listas}
\begin{flushright}
  \footnotesize
  \textit{\textbf{6: Estruturas de dados básicas}}\\
  \textit{c06-listas-operacoes-02}\\

\end{flushright}
Como \haskellinline{String} é uma lista de caracteres, a expressão abaixo pode ser avaliada pelas mesmas funções de listas.
\begin{minted}{haskell}
drop 4 "OuroPreto"
\end{minted}
Qual é o resultado?

\begin{enumerate}
  \item A expressão resulta em erro de tipo.

  \item \haskellinline{ "Ouro" }

  \item \haskellinline{ "reto" }

  \item \haskellinline{ "uroPreto" }

  \item \haskellinline{ "Preto" }


\end{enumerate}

\questionTitle{32}{Inferência da assinatura de areaCirculo}
\begin{flushright}
  \footnotesize
  \textit{\textbf{7: Polimorfismo}}\\
  \textit{cap07-inferencia-area-circulo}\\

\end{flushright}
Considere a definição:

\begin{minted}{haskell}
areaCirculo r = pi * r^2
\end{minted}

Qual é a assinatura de tipo mais geral?

\begin{enumerate}
  \item \mintinline{haskell}{areaCirculo :: Double -> Int}
  \item \mintinline{haskell}{areaCirculo :: Num a => a -> a}
  \item \mintinline{haskell}{areaCirculo :: Ord a => a -> Bool}
  \item \mintinline{haskell}{areaCirculo :: Floating a => a -> a}
  \item \mintinline{haskell}{areaCirculo :: Integral a => a -> a}

\end{enumerate}

\questionTitle{33}{Modelagem de inventário}
\begin{flushright}
  \footnotesize
  \textit{\textbf{6: Estruturas de dados básicas}}\\
  \textit{c06-modelagem-02}\\

\end{flushright}
Você quer representar um inventário como uma lista de produtos. Cada produto possui nome (\haskellinline{String}), código (\haskellinline{Int}) e preço opcional (\haskellinline{Maybe Double}). Qual tipo representa corretamente esse inventário?

\begin{enumerate}
  \item \haskellinline{[(String, Int, Maybe Double)]}

  \item \haskellinline{[String, Int, Maybe Double]}

  \item \haskellinline{([String], [Int], [Maybe Double])}

  \item \haskellinline{Maybe (String, Int, Double)}

  \item \haskellinline{(String, Int, Maybe Double)}


\end{enumerate}

\questionTitle{34}{Análise de assinatura de tipo}
\begin{flushright}
  \footnotesize
  \textit{\textbf{4: Tipos de Dados}}\\
  \textit{c04-assinatura-analise-01}\\

\end{flushright}
Qual das seguintes definições de função corresponde à assinatura de tipo \haskellinline{String -> Int -> Bool}?

\begin{enumerate}
  \item \begin{haskellcode}
verifica b s = not b && null s
\end{haskellcode}

  \item \begin{haskellcode}
verifica n s = n > 0 && s == "OK"
\end{haskellcode}

  \item \begin{haskellcode}
verifica s = s == "OK"
\end{haskellcode}

  \item \begin{haskellcode}
verifica s n = length s > n
\end{haskellcode}

  \item \begin{haskellcode}
verifica s n = s ++ show n
\end{haskellcode}


\end{enumerate}

\questionTitle{35}{Caso base inalcançável}
\begin{flushright}
  \footnotesize
  \textit{\textbf{9: Funções recursivas}}\\
  \textit{c09-fundamentos-terminacao-01}\\

\end{flushright}
Considere a definição abaixo:
\begin{haskellcode}
fatorial :: Integer -> Integer
fatorial n
  | n == 0 = 1
  | otherwise = n * fatorial (n - 1)
\end{haskellcode}
O que acontece ao avaliar \haskellinline|fatorial (-1)|?

\begin{enumerate}
  \item Ocorre imediatamente um erro de padrões não exaustivos, antes de qualquer chamada recursiva.

  \item A função retorna \haskellinline|1|, pois \haskellinline|0| é o caso base mais próximo.

  \item A função calcula \haskellinline|fatorial 1|, ignorando automaticamente o sinal negativo.

  \item A computação não alcança o caso base: ela chama \haskellinline|fatorial (-2)|, depois \haskellinline|fatorial (-3)|, e assim por diante.

  \item A função retorna \haskellinline|0|, pois fatoriais de números negativos são tratados como zero em Haskell.


\end{enumerate}

\questionTitle{36}{Soma produzida por leitura até sentinela}
\begin{flushright}
  \footnotesize
  \textit{\textbf{10: Ações de E/S recursivas}}\\
  \textit{c10-rastreamento-01}\\

\end{flushright}
A função abaixo é executada, e o usuário digita \haskellinline|4, -1, 9, 0|, um valor por linha.

\begin{haskellcode}
lerESomar :: IO Integer
lerESomar = do
  n <- readLn
  if n == 0
    then return 0
    else do
      somaResto <- lerESomar
      return (n + somaResto)
\end{haskellcode}

Qual valor é produzido por \haskellinline|lerESomar|?

\begin{enumerate}
  \item 13

  \item 0

  \item 3

  \item 9

  \item 12


\end{enumerate}

\questionTitle{37}{Raiz quadrada inteira}
\begin{flushright}
  \footnotesize
  \textit{\textbf{9: Funções recursivas}}\\
  \textit{c09-projeto-raiz-inteira-01}\\

\end{flushright}
A raiz quadrada inteira de \haskellinline|n| é o maior \haskellinline|r| tal que \haskellinline|r^2 <= n|. Qual estratégia recursiva está alinhada com essa definição?

\begin{enumerate}
  \item Começar em \haskellinline|0| e diminuir o candidato até que ele seja negativo.

  \item Começar com um candidato \haskellinline|i| e diminuí-lo enquanto \haskellinline|i^2 > n|; quando essa condição falhar, \haskellinline|i| é a resposta.

  \item Somar \haskellinline|n| consigo mesmo até ultrapassar \haskellinline|n^2|.

  \item Multiplicar \haskellinline|n| por \haskellinline|2| até que o resultado seja quadrado perfeito.

  \item Usar \haskellinline|sqrt| e depois converter o resultado para inteiro, pois o exercício exige ponto flutuante.


\end{enumerate}

\questionTitle{38}{Sintaxe explícita e regra de layout}
\begin{flushright}
  \footnotesize
  \textit{\textbf{3: Expressões e definições}}\\
  \textit{c03-layout-avaliacao-02}\\

\end{flushright}
Qual trecho é equivalente à expressão com delimitadores explícitos \haskellinline|let { a = 1; b = 2 } in a + b|?

\begin{enumerate}
  \item \begin{haskellcode}
let a = 1
    b = 2
in a + b
\end{haskellcode}

  \item \begin{haskellcode}
let a = 1
  b = 2
in a + b
\end{haskellcode}

  \item \begin{haskellcode}
a = 1
b = 2
in a + b
\end{haskellcode}

  \item \begin{haskellcode}
let a = 1; b = 2; a + b
\end{haskellcode}

  \item \begin{haskellcode}
let a = 1
in b = 2
in a + b
\end{haskellcode}


\end{enumerate}

\questionTitle{39}{Tipo da condição}
\begin{flushright}
  \footnotesize
  \textit{\textbf{5: Expressão condicional}}\\
  \textit{c05-if-regras-02}\\

\end{flushright}
Qual das expressões abaixo falha especificamente porque a condição do \haskellinline|if| não tem tipo \haskellinline|Bool|?

\begin{enumerate}
  \item \haskellinline|if 3 < 5 then 10 else 20|

  \item \haskellinline|if even 4 then 10 else 20|

  \item \haskellinline|if 3 then 10 else 20|

  \item \haskellinline|if True then 10 else 20|

  \item \haskellinline|if null [] then 10 else 20|


\end{enumerate}

\questionTitle{40}{Avaliação com where e dependência local}
\begin{flushright}
  \footnotesize
  \textit{\textbf{3: Expressões e definições}}\\
  \textit{c03-definicoes-where-02}\\

\end{flushright}
Considere a definição:

\begin{haskellcode}
f x = a + b
  where
    a = x + 1
    b = a * 2
\end{haskellcode}

Qual é o valor de \haskellinline|f 3|?

\begin{enumerate}
  \item \haskellinline|12|.

  \item Ocorre erro, pois \haskellinline|b| não pode usar \haskellinline|a| definido no mesmo bloco \haskellinline|where|.

  \item \haskellinline|7|.

  \item \haskellinline|4|.

  \item \haskellinline|8|.


\end{enumerate}

\questionTitle{41}{Acesso composto a pares}
\begin{flushright}
  \footnotesize
  \textit{\textbf{6: Estruturas de dados básicas}}\\
  \textit{c06-tuplas-acesso-01}\\

\end{flushright}
Considere a expressão abaixo.
\begin{minted}{haskell}
fst (snd (0, ("Carla", 42)))
\end{minted}
Qual é o seu resultado?

\begin{enumerate}
  \item \haskellinline{ 0 }

  \item \haskellinline{ 42 }

  \item A expressão resulta em erro de compilação.

  \item \haskellinline{ "Carla" }

  \item \haskellinline{ ("Carla",42) }


\end{enumerate}

\questionTitle{42}{A primeira linguagem funcional prática}
\begin{flushright}
  \footnotesize
  \textit{\textbf{1: Paradigmas de programação}}\\
  \textit{c01-historia-evolucao-02}\\

\end{flushright}
Criada na virada da década de 1950 por John McCarthy no MIT, qual linguagem de programação é historicamente consagrada como o primeiro esforço prático para incorporar o paradigma funcional, introduzindo conceitos inovadores à época, como a recursão nativa e o \emph{Garbage Collection}?

\begin{enumerate}
  \item Miranda

  \item Lisp

  \item Fortran

  \item COBOL

  \item ML (MetaLanguage)


\end{enumerate}

\questionTitle{43}{Refatoração para forma mais concisa}
\begin{flushright}
  \footnotesize
  \textit{\textbf{5: Expressão condicional}}\\
  \textit{c05-if-vs-guardas-02}\\

\end{flushright}
Considere a função:
\begin{haskellcode}
classifica n
  | n > 0     = "Positivo"
  | n < 0     = "Negativo"
  | n == 0    = "Zero"
  | otherwise = "Impossível"
\end{haskellcode}
Qual definição é funcionalmente equivalente, mas mais concisa?

\begin{enumerate}
  \item \begin{haskellcode}
classifica n
  | n > 0     = "Positivo"
  | n < 0     = "Negativo"
  | otherwise = "Zero"
\end{haskellcode}

  \item \begin{haskellcode}
classifica n
  | n > 0     = "Positivo"
  | otherwise = "Negativo"
\end{haskellcode}

  \item \begin{haskellcode}
classifica n
  | n >= 0    = "Positivo"
  | otherwise = "Negativo"
\end{haskellcode}

  \item A função original não pode ser simplificada sem mudar seu comportamento.

  \item \begin{haskellcode}
classifica n
  | n == 0    = "Zero"
  | otherwise = "Positivo ou Negativo"
\end{haskellcode}


\end{enumerate}

\questionTitle{44}{Literais inteiros e fromInteger}
\begin{flushright}
  \footnotesize
  \textit{\textbf{7: Polimorfismo}}\\
  \textit{cap07-literais-frominteger-contexto}\\

\end{flushright}
O literal inteiro \haskellinline{42} pode aparecer em diferentes tipos numéricos. Qual mecanismo explica esse comportamento?
\begin{enumerate}
  \item Todo literal inteiro tem sempre tipo \haskellinline{Int}; outros tipos recebem coerção implícita em tempo de execução.
  \item Literais inteiros são strings que só são convertidas por \haskellinline{read}.
  \item Literais inteiros pertencem diretamente à classe \haskellinline{Integral}, nunca à classe \haskellinline{Num}.
  \item O compilador cria uma subclasse nova para cada literal encontrado no programa.
  \item Literais inteiros são traduzidos por meio de \haskellinline{fromInteger}, resultando inicialmente em um tipo da forma \haskellinline{Num a => a}.

\end{enumerate}

\questionTitle{45}{Identificação de violação de pureza (I/O temporal)}
\begin{flushright}
  \footnotesize
  \textit{\textbf{1: Paradigmas de programação}}\\
  \textit{c01-transparencia-referencial-02}\\

\end{flushright}
Considere uma função \haskellinline|obterHoraAtual| que consulta o relógio interno do sistema operacional e retorna uma representação do momento exato em que foi chamada.

Por que a existência desta função viola a propriedade de pureza em um modelo matemático estrito?

\begin{enumerate}
  \item Porque seu valor de retorno depende de um estado global em mutação constante fora do controle do programa, fazendo com que chamadas com os mesmos argumentos produzam resultados diferentes.

  \item Porque a leitura do relógio do sistema exige chamadas de baixo nível em \emph{Assembly}, que quebram as abstrações do compilador Haskell.

  \item Porque ela consome excessivos ciclos de processamento do \emph{garbage collector}, atrasando a execução do fluxo principal da aplicação.

  \item Ela não viola a pureza; contanto que a função não modifique a hora do sistema (apenas efetue uma leitura), ela continua sendo referencialmente transparente.

  \item Porque o tempo é uma grandeza contínua e linguagens funcionais só conseguem modelar estruturas de dados e matrizes estritamente discretas.


\end{enumerate}

\questionTitle{46}{Benefício de testar funções puras separadamente}
\begin{flushright}
  \footnotesize
  \textit{\textbf{10: Ações de E/S recursivas}}\\
  \textit{c10-separacao-puro-02}\\

\end{flushright}
Ao projetar um programa interativo em Haskell, por que é desejável separar a coleta de dados em \haskellinline|IO| do processamento puro?

\begin{enumerate}
  \item Porque o compilador exige que todo bloco \haskellinline|do| contenha exatamente uma função pura.

  \item Porque funções puras imprimem seus resultados automaticamente no GHCi e no programa compilado.

  \item Porque funções puras são as únicas que podem chamar \haskellinline|readLn| com segurança.

  \item Porque ações de E/S não podem ter assinaturas de tipo explícitas.

  \item Porque funções puras podem ser testadas e reutilizadas sem depender de entrada e saída reais.


\end{enumerate}

\questionTitle{47}{Avaliação com cálculo local}
\begin{flushright}
  \footnotesize
  \textit{\textbf{5: Expressão condicional}}\\
  \textit{c05-where-avaliacao-02}\\

\end{flushright}
Considere a função:
\begin{haskellcode}
analisa y
  | x > 10    = "Alto"
  | otherwise = "Baixo"
  where
    x = y * y + 1
\end{haskellcode}
Qual é o resultado de \haskellinline|analisa 3|?

\begin{enumerate}
  \item \haskellinline|"Alto"|

  \item \haskellinline|10|

  \item \haskellinline|3|

  \item Ocorre um erro de tipo em \haskellinline|x > 10|.

  \item \haskellinline|"Baixo"|


\end{enumerate}

\questionTitle{48}{Rastreamento do fatorial com acumulador}
\begin{flushright}
  \footnotesize
  \textit{\textbf{9: Funções recursivas}}\\
  \textit{c09-cauda-fat-01}\\

\end{flushright}
Considere:
\begin{haskellcode}
fat' n x
  | n == 0 = x
  | n > 0  = fat' (n - 1) (n * x)
\end{haskellcode}
Qual é a próxima chamada feita por \haskellinline|fat' 4 1|?

\begin{enumerate}
  \item \haskellinline|fat' 3 4|

  \item \haskellinline|fat' 0 24|

  \item \haskellinline|fat' 3 1|

  \item \haskellinline|fat' 4 4|

  \item \haskellinline|fat' 4 3|


\end{enumerate}

\questionTitle{49}{Vantagem de tipos sinônimos}
\begin{flushright}
  \footnotesize
  \textit{\textbf{4: Tipos de Dados}}\\
  \textit{c04-sinonimos-vantagem-01}\\

\end{flushright}
Qual é o principal benefício de usar um tipo sinônimo, como \haskellinline{type Idade = Int}?

\begin{enumerate}
  \item Melhorar o desempenho do programa, pois sinônimos são mais eficientes.

  \item Permitir a definição de tipos recursivos.

  \item É a única maneira de usar \haskellinline{Int} em assinaturas de função.

  \item Criar um tipo completamente novo que é incompatível com \haskellinline{Int}, aumentando a segurança.

  \item Aumentar a legibilidade e a expressividade do código, tornando as assinaturas mais auto-documentadas.


\end{enumerate}

\questionTitle{50}{Bounded e tipos com limites}
\begin{flushright}
  \footnotesize
  \textit{\textbf{7: Polimorfismo}}\\
  \textit{cap07-hierarquia-bounded-integer}\\

\end{flushright}
Qual alternativa explica corretamente por que \haskellinline{Int} é instância de \haskellinline{Bounded}, mas \haskellinline{Integer} normalmente não é?
\begin{enumerate}
  \item \haskellinline{Bounded} só se aplica a tipos booleanos.
  \item \haskellinline{Integer} não é numérico, por isso não pode ter limites.
  \item \haskellinline{Int} tem limites mínimo e máximo dependentes da implementação; \haskellinline{Integer} representa inteiros de precisão arbitrária e não possui limite superior ou inferior fixo.
  \item \haskellinline{Integer} não é instância de \haskellinline{Bounded} porque não é ordenável.
  \item \haskellinline{Int} é instância de \haskellinline{Bounded} porque possui valores fracionários.

\end{enumerate}

\questionTitle{51}{Composição de funções paramétricas em listas}
\begin{flushright}
  \footnotesize
  \textit{\textbf{7: Polimorfismo}}\\
  \textit{cap07-param-head-tail-tipo}\\

\end{flushright}
Qual é a assinatura de tipo mais geral da função abaixo?

\begin{minted}{haskell}
f xs = head (tail xs)
\end{minted}

\begin{enumerate}
  \item \mintinline{haskell}{f :: [a] -> [a]}
  \item \mintinline{haskell}{f :: Ord a => [a] -> a}
  \item \mintinline{haskell}{f :: [a] -> a}
  \item \mintinline{haskell}{f :: Num a => [a] -> a}
  \item \mintinline{haskell}{f :: [[a]] -> [a]}

\end{enumerate}

\questionTitle{52}{Papel de return em ações de E/S recursivas}
\begin{flushright}
  \footnotesize
  \textit{\textbf{10: Ações de E/S recursivas}}\\
  \textit{c10-return-tipos-01}\\

\end{flushright}
Em Haskell, considere a especialização da função \haskellinline|return| usada neste capítulo:

\begin{haskellcode}
return :: a -> IO a
\end{haskellcode}

Qual alternativa descreve corretamente o papel de \haskellinline|return| em uma ação de E/S recursiva?

\begin{enumerate}
  \item Ela imprime automaticamente o valor produzido pela ação de E/S.

  \item Ela interrompe imediatamente a execução da ação atual, como o comando \haskellinline|return| em C, Java ou Python.

  \item Ela reinicia a ação recursiva, funcionando como um comando \haskellinline|continue|.

  \item Ela transforma um valor puro em uma ação \haskellinline|IO| que, ao ser executada, produz esse valor sem realizar interação externa.

  \item Ela lê um novo valor da entrada padrão e o devolve ao próximo passo da recursão.


\end{enumerate}

\questionTitle{53}{Precedência da aplicação sobre operadores}
\begin{flushright}
  \footnotesize
  \textit{\textbf{3: Expressões e definições}}\\
  \textit{c03-precedencia-associatividade-02}\\

\end{flushright}
A expressão \haskellinline|sqrt 16 + 3 * 2| ilustra a precedência da aplicação de função em relação aos operadores infixos. Qual alternativa descreve corretamente sua avaliação?

\begin{enumerate}
  \item Ela é inválida porque operadores aritméticos não podem ser misturados com aplicação prefixa.

  \item Ela é lida como \haskellinline|(sqrt 16) + (3 * 2)|, resultando em \haskellinline|10.0|.

  \item Ela é lida como \haskellinline|sqrt (16 + 3 * 2)|, resultando em aproximadamente \haskellinline|4.6904|.

  \item Ela é lida como \haskellinline|sqrt (16 + 3) * 2|, resultando em aproximadamente \haskellinline|8.7178|.

  \item Ela é inválida porque \haskellinline|sqrt| exige parênteses em volta de \haskellinline|16|.


\end{enumerate}

\questionTitle{54}{Rastreamento do fatorial}
\begin{flushright}
  \footnotesize
  \textit{\textbf{9: Funções recursivas}}\\
  \textit{c09-rastreamento-fatorial-01}\\

\end{flushright}
Considere:
\begin{haskellcode}
fatorial :: Integer -> Integer
fatorial n
  | n == 0 = 1
  | n > 0  = n * fatorial (n - 1)
\end{haskellcode}
Qual é o valor de \haskellinline|fatorial 6|?

\begin{enumerate}
  \item \haskellinline|5040|

  \item \haskellinline|36|

  \item \haskellinline|360|

  \item \haskellinline|120|

  \item \haskellinline|720|


\end{enumerate}

\questionTitle{55}{Escopo limitado à equação}
\begin{flushright}
  \footnotesize
  \textit{\textbf{5: Expressão condicional}}\\
  \textit{c05-where-escopo-02}\\

\end{flushright}
Considere a definição:
\begin{haskellcode}
f x
  | y > 0     = y
  | otherwise = 0
  where y = x - 1

f x = y + 10
\end{haskellcode}
Qual afirmação está correta?

\begin{enumerate}
  \item \haskellinline|where| sempre define variáveis globais no módulo.

  \item A segunda equação não pode usar \haskellinline|y|, pois esse nome vale apenas na primeira equação.

  \item Ambas as equações compartilham \haskellinline|y|, mas apenas se houver uma única assinatura de tipo.

  \item A segunda equação pode usar \haskellinline|y|, mas somente se ela também tiver guardas.

  \item A segunda equação pode usar \haskellinline|y| normalmente, pois \haskellinline|where| tem escopo em toda a função.


\end{enumerate}

\questionTitle{56}{Composição com fromMaybe}
\begin{flushright}
  \footnotesize
  \textit{\textbf{6: Estruturas de dados básicas}}\\
  \textit{c06-maybe-operacoes-01}\\

\end{flushright}
Qual é o resultado da expressão abaixo?
\begin{minted}{haskell}
fromMaybe "pendente" (if True then Just "enviado" else Nothing)
\end{minted}

\begin{enumerate}
  \item \haskellinline{ "pendente" }

  \item \haskellinline{ Nothing }

  \item \haskellinline{ "enviado" }

  \item \haskellinline{ Just "enviado" }

  \item Ocorre um erro de tipos incompatíveis.


\end{enumerate}

\questionTitle{57}{Análise declarativa da abstração do Quicksort}
\begin{flushright}
  \footnotesize
  \textit{\textbf{1: Paradigmas de programação}}\\
  \textit{c01-mecanismo-reducao-02}\\

\end{flushright}
Considere a cláusula principal da implementação declarativa do \emph{Quicksort} em Haskell:

\begin{haskellcode}
qsort (pivo:resto) = qsort (filter (<= pivo) resto) ++ 
                     [pivo] ++ 
                     qsort (filter (> pivo) resto)
\end{haskellcode}

Ao comparar este trecho com a tradicional implementação iterativa em C, qual das seguintes afirmações avalia corretamente a diferença de abstração utilizada na partição dos elementos?

\begin{enumerate}
  \item A versão funcional abstrai a manipulação de índices de memória e laços \cinline|while|, descrevendo a lista particionada como o resultado da aplicação de funções de filtragem baseadas em propriedades lógicas da lista.

  \item A versão funcional falha em implementar a estratégia de \emph{divisão e conquista}, pois a linguagem Haskell não suporta chamadas de recursão de cauda nativa como a linguagem C.

  \item A versão funcional oculta a passagem de ponteiros de memória nos bastidores do operador \haskellinline|:|, realizando uma mutação de array idêntica à troca de variáveis (\cinline|swap|) do algoritmo em C.

  \item A versão funcional requer que a lista de entrada seja previamente ordenada, caso contrário a função \haskellinline|filter| lançará uma exceção de limite de tempo de execução.

  \item A versão funcional é inerentemente mais ineficiente pois exige que o compilador traduza o operador de concatenação (\haskellinline|++|) em loops duplamente encadeados estritos.


\end{enumerate}

\questionTitle{58}{Servidor de Linguagem Haskell (HLS)}
\begin{flushright}
  \footnotesize
  \textit{\textbf{2: Ambiente de desenvolvimento}}\\
  \textit{c02-ecossistema-ferramentas-02}\\

\end{flushright}
Ao configurar o editor Visual Studio Code (VS Code) para desenvolvimento Haskell, instala-se o \emph{Haskell Language Server} (HLS). Qual benefício arquitetural essa ferramenta adiciona ao editor?

\begin{enumerate}
  \item Altera as fontes tipográficas e o esquema de cores do editor para padronizar a exibição visual de mônadas e funtores.

  \item Obriga o VS Code a rodar em um servidor web remoto, permitindo que múltiplos alunos programem no mesmo arquivo simultaneamente.

  \item Compila o código-fonte em segundo plano para enviar o binário final diretamente a plataformas de hospedagem em nuvem.

  \item Fornece funcionalidades de \emph{IDE} inteligentes, como autocompletar, verificação de erros em tempo real e exibição de tipos ao passar o cursor sobre funções.

  \item Sincroniza os arquivos locais automaticamente com as respostas exigidas no Moodle por meio de \emph{commits} contínuos.


\end{enumerate}

\questionTitle{59}{Prompt sem quebra de linha}
\begin{flushright}
  \footnotesize
  \textit{\textbf{8: Programas interativos}}\\
  \textit{c08-buffering-01}\\

\end{flushright}
Um programa usa \haskellinline|putStr| para imprimir um \emph{prompt} e, em seguida, chama \haskellinline|getLine|. Em alguns ambientes, o \emph{prompt} pode não aparecer antes da digitação. Qual é uma solução adequada?

\begin{enumerate}
  \item Trocar \haskellinline|getLine| por \haskellinline|readLn|, pois \haskellinline|readLn| sempre descarrega a saída padrão.

  \item Reordenar o programa colocando \haskellinline|getLine| antes de \haskellinline|putStr|.

  \item Forçar a descarga do \emph{buffer} com \haskellinline|hFlush stdout| logo após \haskellinline|putStr|, ou configurar \haskellinline|stdout| com \haskellinline|NoBuffering|.

  \item Trocar \haskellinline|putStr| por \haskellinline|print|, pois \haskellinline|print| não usa \emph{buffer}.

  \item Adicionar \haskellinline|let| antes de \haskellinline|getLine| para forçar a ordem de execução.


\end{enumerate}

\questionTitle{60}{Comparação entre polimorfismo paramétrico e ad hoc}
\begin{flushright}
  \footnotesize
  \textit{\textbf{7: Polimorfismo}}\\
  \textit{cap07-comparacao-param-adhoc}\\

\end{flushright}
Qual alternativa distingue melhor polimorfismo paramétrico de polimorfismo \emph{ad hoc} em Haskell?
\begin{enumerate}
  \item No polimorfismo paramétrico, a função opera uniformemente para todos os tipos; no polimorfismo \emph{ad hoc}, o comportamento depende de uma instância de classe de tipo.
  \item Polimorfismo paramétrico ocorre apenas em POO; polimorfismo \emph{ad hoc} ocorre apenas em Haskell.
  \item Polimorfismo \emph{ad hoc} não é verificado pelo sistema de tipos.
  \item No polimorfismo paramétrico, cada tipo precisa de uma implementação diferente; no \emph{ad hoc}, há sempre uma única implementação uniforme.
  \item Polimorfismo paramétrico é sinônimo de conversão implícita entre tipos numéricos.

\end{enumerate}

\questionTitle{61}{Análise de erro em length xs / 2}
\begin{flushright}
  \footnotesize
  \textit{\textbf{7: Polimorfismo}}\\
  \textit{cap07-expressao-length-divisao-erro}\\

\end{flushright}
Considere a expressão:

\begin{minted}{haskell}
length [10,20,30] / 2
\end{minted}

Qual é a análise correta?

\begin{enumerate}
  \item A expressão tem erro de tipo, pois \haskellinline{length} retorna \haskellinline{Int}, mas \haskellinline{(/)} exige um tipo \haskellinline{Fractional}.
  \item A expressão avalia para \haskellinline{1.5}, pois Haskell converte \haskellinline{Int} para \haskellinline{Double} automaticamente.
  \item A expressão avalia para \haskellinline{1}, pois \haskellinline{(/)} faz divisão inteira quando os argumentos são inteiros.
  \item A expressão só falha porque a lista tem mais de dois elementos.
  \item A expressão tem tipo \haskellinline{Integral a => a}.

\end{enumerate}

\questionTitle{62}{Apropriação do tipo Integer}
\begin{flushright}
  \footnotesize
  \textit{\textbf{4: Tipos de Dados}}\\
  \textit{c04-tipos-int-integer-01}\\

\end{flushright}
Para qual dos seguintes valores o tipo \haskellinline{Integer} é mais apropriado que \haskellinline{Int}?

\begin{enumerate}
  \item O resultado do cálculo de \haskellinline{factorial 100}.

  \item Um contador em um laço que vai de 1 a 10.

  \item O número de alunos em uma turma.

  \item A idade de uma pessoa.

  \item O número de dias em um mês.


\end{enumerate}

\questionTitle{63}{A ação readLn}
\begin{flushright}
  \footnotesize
  \textit{\textbf{8: Programas interativos}}\\
  \textit{c08-conversao-validacao-01}\\

\end{flushright}
O que a ação \haskellinline|readLn :: Read a => IO a| faz quando é executada?

\begin{enumerate}
  \item Ela lê um único caractere da entrada padrão.

  \item Ela lê uma linha da entrada padrão e tenta convertê-la para um valor de algum tipo pertencente à classe \haskellinline|Read|.

  \item Ela imprime a representação textual de um valor, como \haskellinline|print|.

  \item Ela apenas lê uma linha de texto, sendo idêntica a \haskellinline|getLine|.

  \item Ela valida a entrada automaticamente e, em caso de erro, pede uma nova tentativa.


\end{enumerate}

\questionTitle{64}{Aninhamento com três casos}
\begin{flushright}
  \footnotesize
  \textit{\textbf{5: Expressão condicional}}\\
  \textit{c05-if-aninhamento-02}\\

\end{flushright}
Avalie a função abaixo na entrada \haskellinline|f 5|:
\begin{haskellcode}
f x = if x > 0
      then if x > 10
           then 1
           else 2
      else 3
\end{haskellcode}

\begin{enumerate}
  \item Ocorre um erro de sintaxe por falta de parênteses.

  \item Ocorre um erro de ambiguidade entre os dois \haskellinline|else|.

  \item \haskellinline|1|

  \item \haskellinline|2|

  \item \haskellinline|3|


\end{enumerate}

\questionTitle{65}{Shadowing em expressões let aninhadas}
\begin{flushright}
  \footnotesize
  \textit{\textbf{3: Expressões e definições}}\\
  \textit{c03-let-escopo-02}\\

\end{flushright}
Considere a expressão:

\begin{haskellcode}
let x = 10
    y = x + 1
in let x = 3
   in x + y
\end{haskellcode}

Qual é o seu valor?

\begin{enumerate}
  \item \haskellinline|21|, pois o \haskellinline|x| externo continua valendo em todo o corpo interno.

  \item \haskellinline|6|, pois o \haskellinline|x| interno redefine também o valor de \haskellinline|y|.

  \item Ocorre erro de escopo, pois não é permitido reutilizar o nome \haskellinline|x| em um \haskellinline|let| aninhado.

  \item \haskellinline|14|, pois o \haskellinline|x| interno vale \haskellinline|3| no corpo interno, mas \haskellinline|y| foi definido a partir do \haskellinline|x| externo.

  \item \haskellinline|23|, pois os dois valores de \haskellinline|x| são somados a \haskellinline|y|.


\end{enumerate}

\questionTitle{66}{Regra de nomenclatura de tipos}
\begin{flushright}
  \footnotesize
  \textit{\textbf{4: Tipos de Dados}}\\
  \textit{c04-conceito-nomes-01}\\

\end{flushright}
Qual é a regra sintática obrigatória para nomes de tipos em Haskell?

\begin{enumerate}
  \item Devem começar com uma letra minúscula.

  \item Não podem conter caracteres especiais.

  \item Devem ser escritos inteiramente em maiúsculas.

  \item Devem conter pelo menos um dígito.

  \item Devem começar com uma letra maiúscula.


\end{enumerate}

\questionTitle{67}{O uso exploratório do :help}
\begin{flushright}
  \footnotesize
  \textit{\textbf{2: Ambiente de desenvolvimento}}\\
  \textit{c02-ghci-inspecao-02}\\

\end{flushright}
Qual das seguintes opções descreve a função do comando \texttt{:help} (ou \texttt{:?}) no ambiente interativo GHCi?

\begin{enumerate}
  \item Ele aciona o depurador passo a passo (\emph{debugger}) para ajudar a encontrar falhas lógicas em tempo de execução.

  \item Ele analisa o código que está sendo digitado e sugere o preenchimento automático das funções em tempo real.

  \item Ele exibe um manual interno listando todos os comandos especiais do GHCi e suas respectivas funções.

  \item Ele formata o arquivo atual para corrigir falhas de indentação antes do carregamento.

  \item Ele conecta o terminal à internet para buscar soluções de erros de compilação no fórum oficial da linguagem.


\end{enumerate}

\questionTitle{68}{Primeira guarda verdadeira vence}
\begin{flushright}
  \footnotesize
  \textit{\textbf{5: Expressão condicional}}\\
  \textit{c05-guardas-avaliacao-02}\\

\end{flushright}
Considere a função:
\begin{haskellcode}
f x
  | x >= 0    = 10
  | x > 3     = 20
  | otherwise = 30
\end{haskellcode}
Qual é o valor de \haskellinline|f 4|?

\begin{enumerate}
  \item \haskellinline|10|

  \item \haskellinline|30|

  \item \haskellinline|20|

  \item Ocorre um erro, pois duas guardas são verdadeiras.

  \item O compilador rejeita a definição pela sobreposição das guardas.


\end{enumerate}

\questionTitle{69}{Homogeneidade de listas}
\begin{flushright}
  \footnotesize
  \textit{\textbf{4: Tipos de Dados}}\\
  \textit{c04-checagem-listas-01}\\

\end{flushright}
Em Haskell, se um programador tentar avaliar a expressão \haskellinline{[1, True, "texto"]}, o que acontecerá?

\begin{enumerate}
  \item Ocorrerá um erro apenas em tempo de execução, ao tentar usar os elementos.

  \item Ocorrerá um erro em tempo de compilação, pois todos os elementos de uma lista devem ter o mesmo tipo.

  \item O compilador criará uma tupla automaticamente.

  \item A expressão será aceita e o tipo inferido será \haskellinline{Any}.

  \item Os valores serão implicitamente convertidos para \haskellinline{String}.


\end{enumerate}


\end{multicols}
\makeanswersheet{UFOP}{BCC222: Programação Funcional}{José Romildo}{2026/1}{Primeira Prova}{03}{69}{25}{7}
\gabarito{03}{E,E,A,D,A,A,E,A,A,A,B,B,B,D,A,A,E,A,C,E,E,D,B,B,D,C,B,D,C,B,E,D,A,D,D,E,B,A,C,A,D,B,A,E,A,E,E,A,E,C,C,D,B,E,B,C,A,D,C,A,A,A,B,D,D,E,C,A,B}
\end{document}
